% Aqui começa o capítulo de Introdução.
\chapter{Introdução}\label{chp:Introducao}

Projetos de tecnologia aplicados à área social apresentam desafios particulares relacionados à complexidade dos problemas atendidos, à diversidade de stakeholders envolvidos e à necessidade de rastreabilidade e confiabilidade das informações. Em muitos contextos institucionais, a gestão de casos sociais ainda depende de registros fragmentados, anotações manuais e sistemas pouco integrados, o que dificulta o acompanhamento dos atendimentos, a análise territorial e a tomada de decisão baseada em dados.

Nesse cenário, a adoção de soluções digitais voltadas à gestão estruturada de casos sociais torna-se um fator crítico para o aumento da eficiência operacional, da transparência e da qualidade dos serviços prestados à população. Entretanto, o desenvolvimento desse tipo de sistema ocorre em ambientes marcados por elevada incerteza, mudanças frequentes de requisitos e forte influência de fatores humanos e organizacionais, o que impõe limites às abordagens tradicionais de gestão de projetos baseadas em planejamento rígido.

Este trabalho propõe a simulação do desenvolvimento de um sistema de gestão de casos sociais, concebido como uma aplicação digital com funcionalidades de cadastro e atualização de casos (CRUD), categorização por meio de tags temáticas, registro de apontamentos evolutivos, visualização georreferenciada dos atendimentos em mapa e geração de resumos analíticos para apoio à gestão. O sistema tem como objetivo apoiar equipes técnicas e gestores na organização do trabalho, no acompanhamento longitudinal dos casos e na produção de informações consolidadas para planejamento e avaliação de políticas públicas.

Do ponto de vista da gestão do projeto, a solução é planejada e analisada à luz dos princípios e domínios de desempenho definidos pelo \textcite{PMBOK7}, integrados à utilização de métodos ágeis, com ênfase no framework Scrum. Essa combinação permite explorar uma abordagem orientada à entrega contínua de valor, à adaptação às mudanças e à colaboração entre os envolvidos, sem abrir mão de previsibilidade, controle e alinhamento estratégico.

O projeto é estruturado de forma incremental e iterativa, com ciclos curtos de desenvolvimento, inspeção contínua dos resultados e adaptação do planejamento conforme o aprendizado obtido ao longo da execução. Dessa forma, o trabalho busca simular condições próximas às encontradas no mercado de trabalho, evidenciando como a gestão ágil de projetos pode ser aplicada de maneira consistente e alinhada às diretrizes contemporâneas do PMBOK 7ª edição.

A partir dessa contextualização, os capítulos seguintes apresentam o referencial teórico adotado, a aplicação dos princípios e domínios de desempenho do PMBOK ao projeto, o tailoring realizado, o planejamento ágil com backlog e roadmap, bem como a análise de riscos, incertezas e considerações finais.
